
\section{Background} \label{background}

This section of the report covers relevant background knowledge necessary to understand the developed pose tracking system. A high level overview of a number of computer vision related topics is given.


\subsection{Object Detection and Pose Estimation} \label{detection}

In object detection the goal is to localize all object, the number of which is not known a priori, in images \parencite{Lanz2026Notes}. Like many computer vision recognition tasks, the field of object detection has evolved significantly with the introduction of deep learning \parencite{Lanz2026Notes}. Deep learning-based object detection has evolved through two primary paradigms. Anchor-based models, such as R-CNN or YOLO, rely on predefined bounding box priors of various scales and predicting offsets from these anchors. In contrast, anchor-free models predict bounding boxes directly from point features or keypoint combinations.

In two stage detector models such as R-CNN, there are two separate inference stages \parencite{bharati2019deep}. The first stage generates a large number of candidate regions and the second stage then classifies these regions. Fast R-CNN \parencite{girshick2015fast} and Faster R-CNN \parencite{ren2015faster} improve inference speed over the original R-CNN by using a shared feature backbone for both stages and by using a Region Proposal Network to create candidate regions.

In contrast to these two-stage detectors, the YOLO model used in this project is a single stage detector \parencite{jiang2022review}. It works by dividing the input image into multiple grids of different resolutions. Then a CNN based architecture is employed to extract features for each of the grid positions. Finally, one or more prediction head are applied to the resulting feature maps to simultaneously predict both object presence and bound box offsets from grid location defined anchors. The main advantage of single-stage approaches is that they are usually faster than two-stage approaches, although they often perform worse in practice \parencite{Lanz2026Notes}.

Once an object is detected, one can append extra heads to the feature maps to also predict other properties about the objects. An example of this approach is the Mask R-CNN architecture \parencite{he2017mask}, which predicts instance masks for each detected object \parencite{Lanz2026Notes}. However, this approach can also be applied to Human Pose Estimation, where the additional output head of the network is instead used to predict keypoint locations \parencite{maji2022yolo}. This approach is taken by the YOLO26 models used in this project \parencite{ultralyticsYOLO,sapkota2025yolo26,sapkota2025ultralytics}. Furthermore, as part of this work an additional custom head is trained for the YOLO26 model as explained in \cref{yolo-head}.


\subsection{Object Tracking} \label{tracking}

Detecting objects or estimating poses in individual frames is insufficient for many applications \parencite{Lanz2026Notes}. Instead, what is often required is to associate detections with the physical objects they belong to and model how those evolve over time. This problem is called object tracking. The approach taken in this project is tracking-by-detection. This method can be divided into three essential steps, detection, association, and filtering. In detection, candidate objects are detected in individual frames. Then, these detections are assigned to concrete objects in the association step. Finally, filtering is used to determine the most probable state of the object given the associated detections. Filtering is responsible for refining the noise inherently present in the detection using probabilistic models and preceding observations. This subsection will focus on the association and filtering steps, while the detection step has already been discussed in \cref{detection},

In many application multiple targets must be tracked at the same time. To understand which detection belongs to which underlying object, association must be performed \parencite{Lanz2026Notes}. A typical way to proceed it to predict object from the previous frame and to use those predictions to associate detections to the objects. This association can be performed using an association measure, which is a measure indicating the distance or similarity between a predicted object track and a detection. A naive approach, and also the one applied in this project, is to measure the overlap of the prediction and the detection \parencite{bewley2016simple}. Alternative approaches can include color histograms or optical flow and measure agreement \parencite{Lanz2026Notes}. It is also possible to used metric learning, where a deep learning model is used to extract features for computing the association measure.

After an association measure has been established, it is still necessary to actually perform the association \parencite{Lanz2026Notes}. A naive approach can be to consider all detections within some gating area of the prediction and then associate the closest of these detections to the prediction. The distance in this context can also consider the probabilistic model of the prediction to find most probable detections. However, this nearest neighbor approach can easily fail due to ambiguities. A better approach is to use a bipartite graph matching based framework. In this framework, the problem is modeled as a bipartite graph from predictions to detections, with edges weighted by the association metric. This can be formulated as the following integer linear programming problem:
\begin{align*}
  & \operatorname*{argmax}_{\mathbf{Z}} \sum_{i=1}^{M} \sum_{j=1}^{N} s_{ij} z_{ij} \\
  & s.t. \quad \forall j: \sum_{i=1}^{M} z_{ij} = 1 \quad \forall i: \sum_{j=1}^{N} z_{ij} = 1 \quad z_{ij} \in \{0, 1\}
\end{align*}
where $s_{ij}$ and $z_{ij}$ are respectively the association metric and association between prediction $i$ and detection $j$. This problem can then be solved either using a greedy algorithm, or as has been done for this project, using globally optimal algorithms such as the Hungarian method \parencite{kuhn1955hungarian}. Finally, another possible approach to association is to used advanced graph-based deep learning networks.

After detections and objects have been associated, it is necessary to perform the filtering step \parencite{Lanz2026Notes}. It is crucial to understand that detections differ from the underlying state that we try to understand. Aside from being noisy, detections might also be missing some information of interest in the object. For example, object detections usually do not include velocity, and still we would like to infer it from the detections. When represented as using a Markov model, the state of a target at any given time step depends only on that targets state in the previous time step. Further, the observation or detection at any time step depends only on the targets current state. One way to evaluate this process is using recursive Bayesian estimation, also called Bayesian filtering.

Bayesian filtering is applied in two phases. First, the prediction step is used to model the target states distribution using only the past observations. Due to the conditional independence assumption of the hidden Markov model, this probability distribution of the hidden state can be computed as followed:
\begin{align*}
  p(\mathbf{x}_t|\mathbf{z}_{1:t-1}) &= \int p(\mathbf{x}_t, \mathbf{x}_{t-1} | \mathbf{z}_{1:t-1}) d\mathbf{x}_{t-1} \\
  &= \int p(\mathbf{x}_t | \mathbf{x}_{t-1}, \mathbf{z}_{1:t-1}) p(\mathbf{x}_{t-1} | \mathbf{z}_{1:t-1}) d\mathbf{x}_{t-1} \\
  &= \int p(\mathbf{x}_t | \mathbf{x}_{t-1}) p(\mathbf{x}_{t-1} | \mathbf{z}_{1:t-1}) d\mathbf{x}_{t-1}
\end{align*}
where $p(\mathbf{x}_t | \mathbf{x}_{t-1})$ defines out motion model, and $p(\mathbf{x}_{t-1} | \mathbf{z}_{1:t-1})$ is the posterior distribution after applying all previous observations. After prediction, we want to incorporate the information gained by the new observation. This can be achieved as follows by applying Bayes' rule:
\begin{align*}
  p(\mathbf{x}_t \mid \mathbf{z}_{1:t}) &\propto p(\mathbf{z}_t \mid \mathbf{x}_t) p(\mathbf{x}_t \mid \mathbf{z}_{1:t-1})
\end{align*}
where $p(\mathbf{z}_t \mid \mathbf{x}_t)$ models the observation process.

The equations of the Bayesian filtering algorithm can be implemented by representing the probability distributions at each time step in parametric form, such as in Kalman filtering \parencite{Lanz2026Notes,bar2001estimation}, or in nonparametric form, such as with particle filters. In a linear Kalman filter the motion and observation models are defined as:
\begin{align*}
  \mathbf{x}_t &= \mathbf{F}_t \mathbf{x}_{t-1} + \mathbf{B}_t \mathbf{u}_t + \mathbf{w}_t, \quad \mathbf{w}_t \sim \mathcal{N}(\mathbf{0}, \mathbf{Q}_t) \\
  \mathbf{z}_t &= \mathbf{H}_t\mathbf{x}_t + \mathbf{v}_t, \quad \mathbf{v}_t \sim \mathcal{N}(\mathbf{0}, \mathbf{R}_t)
\end{align*}
The optimal closed-form solution is the standard Kalman Filter, which updates the mean $\mathbf{\bar{x}}_t$ and covariance $\Sigma_t$ of $\mathbf{x}_t$ through linear projection and gain computation.

While the linear Kalman filter is applicable in cases where both the motion and observation model are linear, in our application of multi-view tracking, the measurement model is non-linear because 3D world coordinates must project onto a 2D camera plane. For this reason, this project makes use of extended Kalman filtering. We define the observation function to be $\mathbf{z}_t = h(\mathbf{x}_t) + \mathbf{v}_t$. The extended Kalman filter linearizes this non-linear function using a first-order Taylor expansion around the predicted state estimate $\mathbf{\bar{x}}_t$:
\begin{align*}
h(\mathbf{x}_t) \approx h(\mathbf{\bar{x}}_t) + \mathbf{J}_h (\mathbf{x}_t - \mathbf{\bar{x}}_t)
\end{align*}
where $\mathbf{J}_h$ is the Jacobian matrix evaluated at $\mathbf{\bar{x}}_t$. In the update step the Jacobian matrix is the used in place of the linear observation matrix $\mathbf{H}_t$ to update the covariance and calculate Kalman gain.


\subsection{Projective Geometry} \label{calibration}

The dataset used in this project already contains camera calibration parameters for all camera views. Still, understanding the meaning of these parameters and how they relate to the image formation process is critical for the correct implementation for the triangulation algorithm as well as for modeling the observation process. The mapping of 3D points to 2D pixels is defined mathematically through the pinhole camera model \parencite{Lanz2026Notes,hartley2003multiple}. Unlike a physical pinhole camera, which would invert the image, the mathematical model places the image plane in front of the focal point for convenience. The full perspective projection model linking a homogenous 3D world coordinate $\tilde{x}_w$ to a homogenous 2D sensor pixel $\tilde{x}_s$ is given by:
\begin{align*}
  \tilde{x}_s &= \mathbf{K} [\mathbf{R} | \mathbf{t}] \tilde{x}_w = \mathbf{P} \tilde{x}_w
\end{align*}
where $\mathbf{P}$ is the full $3 \times 4$ projection matrix. It decomposes into intrinsic parameters $\mathbf{K}$ and extrinsic parameters $\mathbf{R} and \mathbf{t}$. For a perspective projection, 3D points $\tilde{x}_s$ in the local camera coordinate system are mapped to the image plane by dividing them by their depth $Z$ component.

The assumption of linear projection made by the above pinhole camera model, where straight lines remain straight, is violated by physical camera lenses, introducing distortions \parencite{Lanz2026Notes}. These must be modeled to accurately project points or undistort pixel coordinates. Two types of distortions are typically modeled. First, radial distortion causes points to bend outward inward. It is modeled relative to the squared radius $r^2 = x^2 + y^2$. Second, tangential distortion arises when the lens is not perfectly parallel to the imaging sensor.


\subsection{Multi-View Triangulation} \label{triangulate}

Even once corresponding points across multiple camera views have been identified, their 3D position $\tilde{x}_w$ must still be reconstructed from the 2D pixel locations. Rays back-projected from pixels will nearly never intersect perfectly in 3D space due to measurement noise. As such, a noise robust algorithm must be used that can recover 3D coordinates using multiple noise 2D projections \parencite{Lanz2026Notes,hartley2003multiple}. If we denote $\tilde{x}_i^s$ the homogeneous 2D observation in the $i$-th camera, and $\mathbf{P}_i$ the projection matrix of that that camera, because $\tilde{x}_i^s$ and the projected point $\mathbf{P}_i \tilde{x}_w$ represent the same point up to scale, their cross product must be zero:
\begin{align*}
  \tilde{x}_i^s \times \mathbf{P}_i \tilde{x}_w &= 0 \\
\end{align*}
This cross product yields a set of three linear equations for each camera view. However, the third is a linear combination of the first two, so it can be discarded. By stacking the equations from $N \ge 2$ camera views, we obtain a homogeneous linear system $\mathbf{A} \tilde{x}_w = 0$. Because $\tilde{x}_w$ is defined only up to scale, we impose the constraint $||\tilde{x}_w|| = 1$. The optimal least-squares solution to this system is found by computing the Singular Value Decomposition of $\mathbf{A}$ and taking the right singular vector associated with the smallest singular value. Alternatively, it is also possible to transform the homogeneous system in to an in-homogeneous system and solve it using standard least squares. This latter approach is taken in the project.
